\chapter{کلیات پروژه}
\label{chap:introduction}

\section{مقدمه}
با زیاد شدن منابع تجدیدپذیر پراکنده و ذخیره‌سازهای کوچک، دیگر نمی‌توان شبکه را فقط به صورت تولید متمرکز و مصرف یک‌طرفه دید. ریزشبکه مجموعه‌ای محلی از تولید، باتری و بار است که می‌تواند به شبکه سراسری وصل بماند یا در شرایط خاص جدا کار کند. برای کار کردن چنین سامانه‌ای باید در هر ساعت تصمیم گرفت برق از کجا بیاید، باتری چه کند و آیا از شبکه خرید شود یا به آن فروخته شود.

این تصمیم هم جنبه پولی دارد و هم جنبه فیزیکی. قیمت خرید و فروش و استهلاک باتری در هزینه دیده می‌شود. از آن طرف، در هر گام باید توان‌ها جمع جبری درستی بسازند و وضعیت شارژ از بازه مجاز بیرون نرود. مدلی که فقط هزینه را نگاه کند ممکن است برنامه‌ای بدهد که روی کاغذ ارزان است اما اجرا نمی‌شود. مدلی که فقط معادله بنویسد و از داده استفاده نکند هم در برابر تغییر بار و تابش انعطاف کمی دارد.

از دید علوم کامپیوتر مسئله نوعی یادگیری با قید است. شبکه عصبی می‌تواند از روی مشاهده‌ها خروجی بسازد، ولی فضای توابع آن خیلی بزرگ است. یک راه محدود کردن این فضا گذاشتن قانون فیزیکی داخل تابع هزینه است. رئیسی و همکاران این کار را برای معادلات دیفرانسیل نشان دادند \pcite{raissi2019pinns}. اولادجو و فولی هم مدیریت یک ریزشبکه متصل به شبکه را با تولید خورشیدی، باتری و تقسیم سود میان مشارکت‌کنندگان بررسی کرده‌اند \pcite{oladejo2019microgrid}. در این پروژه همان ایده هزینه فیزیک‌آگاه روی قیود بهره‌برداری انرژی پیاده شده است؛ لایه نظریه بازی فقط در مرور منابع آمده و در آزمایش عددی نیست.

\section{بیان مسئله}
فرض کنید در ساعت $t$ بار، توان خورشیدی، قیمت و وضعیت شارژ معلوم باشد. می‌خواهیم توان شبکه، شارژ، دشارژ و وضعیت شارژ ساعت بعد را به دست آوریم. اگر آموزش فقط اختلاف با برچسب را کم کند، مدل روی همان برچسب‌ها خوب به نظر می‌رسد و ممکن است توازن توان را خراب کند.

سؤال اصلی این است که چطور همان قوانین را طوری بنویسیم که هنگام آموزش قابل استفاده باشند. اینجا باقیمانده توازن و باقیمانده دینامیک باتری با وزن مشخص به هزینه اضافه می‌شود. معادلات از جنس مکانیک سیالات مقاله رئیسی نیستند؛ معادلات ساده بهره‌برداری‌اند.

سؤال دوم این است که سهم مقاله ریزشبکه تا کجا است. آن کار بیشتر روی تقسیم منصفانه سود و الگوریتم \lr{TLBO} تمرکز دارد \pcite{oladejo2019microgrid}. در این گزارش آن قسمت کنار گذاشته شده تا ادعا از چیزی که پیاده شده بزرگ‌تر نشود.

\section{اهمیت موضوع}
داده واقعی ریزشبکه معمولاً کم است و تهیه آن راحت نیست. در عوض توازن انرژی و حدود باتری از قبل معلوم‌اند و دور ریختن‌شان منطقی به نظر نمی‌رسد. در لایه بهره‌برداری، خروجی خارج از قید عملاً به درد نمی‌خورد؛ بنابراین فقط نگاه کردن به میانگین مربع خطا کافی نیست. بعد از آموزش هم ارزیابی شبکه معمولاً از حل دوباره یک مسئله بهینه‌سازی بزرگ سریع‌تر است.

برای پروژه کارشناسی همین اندازه کافی است که مسئله روشن تعریف شود، پیاده‌سازی قابل اجرا باشد و نتیجه با همان داده گزارش شود. انتظار کشف قانون فیزیکی جدید در این سطح نیست.

\section{اهداف}
هدف کلی، آزمودن یک شبکه فیزیک‌آگاه روی مدیریت انرژی یک ریزشبکه متصل به شبکه است. به‌طور جزئی‌تر:
\begin{enumerate}
  \item جدا کردن ایده \lr{PINN} از مثال‌های معادله دیفرانسیل جزئی مقاله مبنا؛
  \item نوشتن متغیرها و قیود یک مدل ساده ریزشبکه؛
  \item تعریف ورودی، خروجی و هزینه ترکیبی؛
  \item پیاده‌سازی با پایتون و مقایسه با مدل فقط‌داده؛
  \item روشن نگه داشتن این که نظریه بازی در عددها نیست.
\end{enumerate}

\section{پرسش‌ها}
\begin{enumerate}
  \item توازن توان و دینامیک باتری را چطور می‌توان در هزینه نوشت؟
  \item زیاد کردن وزن فیزیک با خطای داده چه نسبتی پیدا می‌کند؟
  \item غیر از \lr{MSE} چه چیزی باید گزارش شود؟
  \item کدام عدد مال این پروژه است و کدام مال مقاله‌های مبنا؟
\end{enumerate}

\section{فرض‌ها}
ریزشبکه متصل به شبکه فرض شده است. منبع تجدیدپذیر فقط خورشید است و از دینامیک مبدل صرف‌نظر شده. باتری با یک رابطه انرژی گسسته و بازده ثابت مدل شده است. گام زمانی یک ساعت است. پایداری گذرا بررسی نشده است. عددهای داخل دو مقاله مبنا به‌عنوان نتیجه این کار تکرار نشده‌اند.

\section{روش کار}
ابتدا چارچوب هزینه فیزیک‌آگاه از \pcite{raissi2019pinns} خوانده شد. بعد متغیرهای توان و باتری از مدل مدیریت انرژی \pcite{oladejo2019microgrid} استخراج گردید. سپس یک مدل کوچک نوشته و روی داده شبیه‌سازی‌شده با بذر ثابت آزمایش شد. جزئیات اجرا در فصل~\ref{chap:implementation} است.

\section{ساختار گزارش}
فصل~\ref{chap:background} مبانی را می‌آورد. فصل~\ref{chap:method} مدل را می‌نویسد. فصل~\ref{chap:implementation} پیاده‌سازی و عددها را گزارش می‌کند. فصل~\ref{chap:conclusion} جمع‌بندی است.
